\documentclass[a4paper,12pt]{letter}

\usepackage{amsmath}
\usepackage{amssymb}
\usepackage{nopageno}
\usepackage{amsmath}
\usepackage{enumitem}
% \usepackage{mathtools}
%\usepackage{eqnarray,amsmath}
\usepackage[a4paper, total={7.5in, 11.25in}]{geometry}

\begin{document}

Name: \hrulefill\ Neptun code: \hrulefill

\bigskip

\begin{center}
(KTXFI2EBNF) Physics II. Exam~2 Solutions\\
January~5,~2026
\end{center}

\bigskip
 
\begin{enumerate}[label=\Roman*.]
   
\item Moving Reference Frames. Special Theory of Relativity

\begin{enumerate}[label=\arabic*.]
\setcounter{enumii}{0} % Számláló beállítása

\item State the Galilean principle of relativity. \hfill (5~points)

\smallskip
\textit{Lecture on October 13, 2025, Slides 8--16.}

\item Draw the experimental setup of the Michelson–Morley interferometer, indicate and name the instruments involved in the experiment. What was the objective of performing this experiment? What was the outcome of the experiment? \hfill (5~points)

\smallskip
\textit{Lecture on October 13, 2025, Slide 17.}


\end{enumerate}

% ------------------------------------------------------------------------------------------------------------------------------------

\item Limits of the Classical Conceptual Framework

\begin{enumerate}[label=\arabic*.]
\setcounter{enumii}{2} % Számláló beállítása

\item What are the Stefan–Boltzmann law and Wien’s displacement law? Describe them in your own words and also give their mathematical expressions. Specify the physical meaning of the symbols appearing in the formulas. \hfill (5~points)

\smallskip
\textit{Lecture on October 6, 2025, Slides 12--14.}
  
\item What is the Compton effect, what did it prove, and how was this demonstrated?  \hfill (5~points)

\textit{Lecture on October 6, 2025, Slides 29--31.}

\item Describe Planck’s hypotheses; that is, explain the significance of Planck’s radiation law.  \hfill (5~points)

\textit{Lecture on October 6, 2025, Slides 18--22.}

\end{enumerate}
  
% ------------------------------------------------------------------------------------------------------------------------------------

\item Elements of Quantum Mechanics. Physics of Condensed Matter. Theory of Electrical Conduction

\begin{enumerate}[label=\arabic*.]
\setcounter{enumii}{5} % Számláló beállítása
  
\item Describe the electrical conduction mechanism of semiconductors using the band theory of solids. What characterizes semiconductors from the point of view of electrical conduction? \hfill (5~points)

\smallskip
\textit{Lecture on October 27, 2025, Slides 23, 28.}

\item What is the tunneling effect? Describe the phenomenon in your own words. \hfill (5~points)

\smallskip
\textit{Lecture on November 10, 2025, Slide 46.}

\end{enumerate}

% ------------------------------------------------------------------------------------------------------------------------------------
  
\item Development of Atomic Theory

\begin{enumerate}[label=\arabic*.]
\setcounter{enumii}{7} % Számláló beállítása

\item Describe the Rutherford scattering experiment. \hfill (5~points)

\smallskip
\textit{Lecture on October 20, 2025, Slides 5--7.}

\item What is the Zeeman effect? Describe the phenomenon. Under what conditions does it occur and how does it manifest itself? \hfill (5~points)

\smallskip
\textit{Lecture on October 20, 2025, Slides 16--17.}

\item Given the element with atomic number 11 in the periodic table, write down its electronic configuration using the standard atomic physics notation. \hfill (5~points)

\smallskip
\textit{Lecture on October 20, 2025, Slides 19--22.}
  
\end{enumerate}

\end{enumerate}

% ------------------------------------------------------------------------------------------------------------------------------------

Grades: 0–25~points: Fail (1), 26–30 points: Pass (2), 31–37 points: Satisfactory (3), 38–45~points: Good~(4), 46–50 points: Excellent (5).

\rightline{Dr.~Gabor FACSKO}
\rightline{\textit{facsko.gabor@uni-obuda.hu}}

\vfill

\end{document}
